% Requires: \usepackage{tikz} \usetikzlibrary{arrows.meta,calc}
% Generated by make_dataset_stats.py -- do not edit by hand.
% Counts come from cleaning_report.json, the language labels and the hallucination-judge summary.
\begin{figure*}[t]
\centering
\definecolor{plblue}{HTML}{DBE8FB}
\definecolor{plbluetext}{HTML}{1F3F8F}
\definecolor{plyellow}{HTML}{FCEFC4}
\definecolor{plyellowtext}{HTML}{8A4B08}
\definecolor{plgreen}{HTML}{EAF8F1}
\definecolor{plgreenline}{HTML}{6FCFAE}
\definecolor{plgreentext}{HTML}{0B5D44}
\definecolor{plpink}{HTML}{F3E8FF}
\definecolor{plpinktext}{HTML}{5B21B6}
\definecolor{plline}{HTML}{B9BFCC}
\definecolor{plred}{HTML}{B91C1C}
\definecolor{plmuted}{HTML}{4B5563}
\begin{tikzpicture}[
  font=\scriptsize,
  stage/.style={draw=#1, rounded corners=4pt, line width=0.6pt, fill=white, text width=2.95cm,
                minimum height=2.1cm, align=left, inner sep=4pt, anchor=north west,
                execute at begin node={\hyphenpenalty=10000\exhyphenpenalty=10000}},
  head/.style={rounded corners=3pt, inner xsep=4pt, inner ysep=2.5pt, anchor=south west,
               font=\footnotesize\bfseries},
  flow/.style={draw=plred, line width=1pt, -{Triangle[length=4.5pt,width=4.5pt]}},
]
\node[stage=plline] (a) at (0,0) {%
  CVEfixes + Big-Vul\\[3pt]
  {\color{plmuted}per vulnerable function:}\\
  code, CVE ID,\\ CVE description};
\node[stage=plline] (b) at (3.6,0) {%
  Llama-3.1-8B-Instruct\\
  {\color{plmuted}input: code + description}\\[3pt]
  \textbf{110{,}817} records\\
  \textbf{540{,}654} Q--A pairs};
\node[stage=plline] (c) at (7.2,0) {%
  Remove minified code, unrelated snippets, non-answers, duplicates; cap 200 per CVE\\[3pt]
  {\color{plmuted}$-$84{,}814 records}};
\node[stage=plgreenline, fill=plgreen] (d) at (10.8,0) {%
  \textbf{26{,}003} records\\
  \textbf{119{,}118} Q--A pairs\\
  6{,}911 CVEs\\
  ten languages};
\node[stage=plline] (e) at (14.4,0) {%
  1{,}000-record sample\\
  three LLM judges\\[3pt]
  \textbf{14.8\%} of pairs hallucinated\\
  {\color{plmuted}95\% CI 13.5--16.2}};
\node[head, fill=plblue, text=plbluetext] at ($(a.north west)+(0.12,-0.02)$) {1 Data Collection};
\node[head, fill=plblue, text=plbluetext] at ($(b.north west)+(0.12,-0.02)$) {2 QA Generation};
\node[head, fill=plyellow, text=plyellowtext] at ($(c.north west)+(0.12,-0.02)$) {3 Rule-based Filtering};
\node[head, fill=plgreen, text=plgreentext] at ($(d.north west)+(0.12,-0.02)$) {4 \dsname};
\node[head, fill=plpink, text=plpinktext] at ($(e.north west)+(0.12,-0.02)$) {5 Hallucination Audit};
\draw[flow] (a.east) -- (b.west);
\draw[flow] (b.east) -- (c.west);
\draw[flow] (c.east) -- (d.west);
\draw[flow, dashed] (d.east) -- (e.west);
\end{tikzpicture}
\caption{Construction of \dsname. An LLM generates question--answer conversations from each vulnerable
function and its CVE description; rule-based filters remove unusable records and pairs
(Section~\ref{sec:filtering}); a random sample of the result is audited by independent LLM judges
(Section~\ref{sec:hallucination}).}
\label{fig:pipeline}
\end{figure*}
